% Part: lambda-calculus
% Chapter: church-rosser
% Section: definitions-and-properties

\documentclass[../../../include/open-logic-section]{subfiles}

\begin{document}

\olfileid{lam}{cr}{dap}

\olsection{التعريف والخصائص}

نعرّف في هذا الفصل خاصية تشيرش--روسر، ثم نبيّن جملة من الخواص المتصلة بها.

\begin{defn}[خاصية تشيرش--روسر، CR] العلاقة~$\xredone$ على الحدود
  ذات \emph{خاصية تشيرش--روسر} إذا كان تحقق
  $M \xredone P$ و$M \xredone Q$ يستلزم وجود حد~$N$ يحقق
  $P \xredone N$ و$Q \xredone N$.
\end{defn}

يمكن أن نتخذ حساب لامبدا نموذجًا للحساب: فتكون الحدود السوية فيه قيمًا،
وتكون علاقة الاختزال قواعد لإجراء الحساب.
ومؤدى خاصية تشيرش--روسر أن تعدد سبل إجراء الحساب لا يوجب تعدد قيمة العبارة.

ومن الجبر الابتدائي مثال يقرب ذلك؛ فحساب
$4 \times (1+2) + 3$ يجري بأكثر من طريق.
نختزله إلى~$4 \times 3+3$ إذا بدأنا باختزال~$1+2$ إلى~$3$،
أو إلى~$4 \times 1+4 \times 2+3$ إذا بدأنا بتوزيع~$4 \times (1+2)$.
ثم تنتهي كلتا العبارتين بمزيد من الاختزال إلى~$12+3$.

إذا جعلنا~$\xredone$ اختزال~$\beta$، لزم من خاصية تشيرش--روسر أن الصورة السوية،
متى وجدت، لا تتعدد. فلنفرض أن~$M$ يختزل إلى~$P$ وإلى~$Q$،
وأن الناتجين سويان. تقتضي الخاصية وجود حد~$N$ يختزل إليه كل من~$P$ و$Q$.
ولما كان~$P$ و$Q$ سويين، لم يكن اختزالهما إلى~$N$ إلا اختزالًا لا يغيّر شيئًا.
فـ$P$ و$Q$ و$N$ متطابقة، وبهذا يسوغ إطلاق \emph{الصورة} السوية للحد.

فإذا عدّ حساب لامبدا نموذجًا للحساب، كانت الصورة السوية بمنزلة النتيجة النهائية
للحساب الذي ابتدأ من الحد. وقد ثبت أن النتيجة النهائية لا تزيد على واحدة،
إن وجدت؛ كما أن حساب~$4 \times (1+2)+3$ لا يعطي إلا~$15$.

\begin{thm} \ollabel{thm:str}
  لتكن العلاقة~$\xredone$ ذات خاصية تشيرش--روسر، ولتكن~$\xred$ أصغر علاقة متعدية
  تشتمل على~$\xredone$. فإن~$\xred$ أيضًا ذات خاصية تشيرش--روسر.
\end{thm}

\begin{proof}
  لنفرض~$M \xred P$ و$M \xred Q$؛ فتوجد المتتاليتان
  \begin{align*}
    M & \xredone P_1 \xredone \dots \xredone P_m \text{ و}\\
    M & \xredone Q_1 \xredone \dots \xredone Q_n\text{،}
  \end{align*}
  حيث~$P_m=P$ و$Q_n=Q$.
  نثبت المطلوب بإنشاء شبكة~$N$ من الحدود، ارتفاعها~$m + 1$ وعرضها~$n + 1$.
  ونرمز بـ$N_{i,j}$ إلى الحد الذي في الصف~$i$ والعمود~$j$.
  
  وننشئ~$N$ بحيث~$N_{i,j} \xredone N_{i+1,j}$ و$N_{i,j} \xredone N_{i,j+1}$، على الوجه الآتي:
  \begin{align*}
    N_{0,0} &= M \\
    N_{i,0} &= P_i && \text{إذا كان } 1 \le i \le m \\
    N_{0,j} &= Q_j && \text{إذا كان } 1 \le j \le n \\
  \intertext{وفي غير ذلك:}
    N_{i,j} &= R
  \end{align*}
  فيه~$R$ حد يحقق~$N_{i-1,j} \xredone R$ و$N_{i,j-1}
  \xredone R$.
  ووجوده مضمون بخاصية تشيرش--روسر للعلاقة~$\xredone$.

  وينتج عن الإنشاء~$N_{m,0} \xredone \dots \xredone N_{m,n}$ و$N_{0,n} \xredone \dots \xredone N_{m,n}$.
  ثم إن~$N_{m,0}$ هو~$P$، و$N_{0,n}$ هو~$Q$.
  فيحصل المطلوب من تعريف~$\xred$.
\end{proof}

\end{document}

